\chapter{Bluebottle Stings}

\section*{Definition and Overview}
A bluebottle sting is a painful skin reaction caused by contact with the tentacles of the bluebottle jellyfish (\textit{Physalia} species), a common marine animal found on Australian beaches and in waters of the Pacific and Atlantic oceans. The bluebottle has a blue, gas-filled float that rests on the water's surface and long trailing tentacles armed with microscopic stinging cells. When a tentacle brushes against skin, these stinging cells discharge venom into the skin, causing a sharp, burning pain and a characteristic red, raised, beaded rash. Most stings are painful but not dangerous, and first aid alone is usually sufficient.

\section*{Epidemiology}
Bluebottle stings are very common in Australia, particularly along the eastern and southern coastlines, and are among the most frequent causes of marine-related beach injury in summer. The jellyfish is blown inshore by onshore winds, especially during the warmer months, and stings occur most often in swimmers, surfers, and people wading in the surf. Children are frequently affected because they may touch the jellyfish on the beach. Most stings are minor and resolve without medical attention, and severe reactions are rare.

\section*{Aetiology and Pathophysiology}
The sting occurs through a direct contact mechanism:

\begin{itemize}
    \item \textbf{Stinging cells (nematocysts):} the tentacles contain millions of microscopic capsules that discharge on contact with skin, injecting venom
    \item \textbf{Venom effects:} bluebottle venom releases substances that irritate nerve endings in the skin, producing intense local pain and inflammation
    \item \textbf{Triggers for discharge:} contact pressure, friction (such as rubbing the skin), and temperature changes can make remaining stinging cells fire
    \item \textbf{Severity factors:} the area of skin touched, the amount of tentacle contact, and the person's individual sensitivity influence how severe the sting feels
\end{itemize}

\section*{Clinical Features}
Symptoms usually begin immediately and can last from minutes to hours:

\begin{itemize}
    \item A sharp, stinging, or burning pain at the site of contact
    \item A red, raised rash that often appears as a string of beads or a whiplike line
    \item Itching and swelling around the affected area
    \item Pain that may spread along the limb
    \item Most symptoms settle within a few hours, though the skin marking may persist for a few days
    \item Severe reactions are uncommon, but general symptoms such as nausea, headache, or feeling unwell may occur with more extensive contact
\end{itemize}

\section*{Differential Diagnosis}
Other causes of pain and skin irritation around the beach include:

\begin{itemize}
    \item Other jellyfish stings, such as those from the box jellyfish (\textit{Chironex fleckeri}), which can be far more dangerous
    \item Irukandji syndrome, caused by very small jellyfish, with delayed severe back pain, sweating, and distress
    \item Stonefish, sea urchin, or stingray injuries causing puncture wounds and pain
    \item An allergic skin reaction to something else in the water or on the skin
    \item Seabather's eruption or other causes of an itchy rash
\end{itemize}

\section*{When to Seek Medical Care}
Most bluebottle stings can be managed with first aid, but seek medical help if:

\begin{itemize}
    \item The pain is severe or does not settle after rinsing and hot water treatment
    \item The rash covers a very large area, or symptoms such as nausea, vomiting, headache, or faintness develop
    \item Stinging cells are felt to be present in the mouth, eyes, or on the face
    \item There are signs of a secondary infection in the wound over the following days
    \item You are unsure whether the sting was from a bluebottle or from a more dangerous jellyfish
\end{itemize}

\textbf{Call 000 (or local emergency services) for:}
\begin{itemize}
    \item Difficulty breathing, chest pain, or collapse
    \item Any sting accompanied by a rapid allergic-type reaction or loss of consciousness
\end{itemize}

\section*{Diagnosis and Investigations}
Diagnosis is usually made on the description and appearance:

\begin{itemize}
    \item \textbf{History:} the timing of the sting, where it occurred, and what the jellyfish looked like are key clues
    \item \textbf{Examination:} the typical beaded or whiplike pattern of the rash supports the diagnosis
    \item \textbf{Identification:} bluebottles are recognised by their blue float and long trailing tentacles
    \item \textbf{Rarely needed:} laboratory tests or imaging are generally not required for bluebottle stings
    \item \textbf{Assessment of severity:} checking for general symptoms and for signs of a more serious jellyfish envenomation
\end{itemize}

\section*{Management}
For a bluebottle sting, first aid is the main treatment:

\begin{itemize}
    \item \textbf{Rinse:} carefully rinse the affected area with seawater to remove any remaining tentacle fragments (fresh water can cause more stinging cells to fire)
    \item \textbf{Remove tentacles:} pick off any visible tentacle pieces with fingers or tweezers -- light touch will not cause stinging from intact bluebottle tentacles
    \item \textbf{Hot water:} immerse the area in hot (but not scalding) water, around 40--45 degrees Celsius, for 20 minutes, or apply hot packs if immersion is not practical
    \item \textbf{Do not:} rub the area, use vinegar, or apply ice directly to a bluebottle sting -- these are not recommended for this species
    \item \textbf{Pain relief:} paracetamol may be used if pain persists after hot water treatment
    \item \textbf{Observation:} most people recover fully with first aid alone; further treatment is rarely needed
\end{itemize}

\section*{Complications}
Complications are uncommon but possible:

\begin{itemize}
    \item Secondary bacterial infection if the skin is scratched or broken
    \item Hyperpigmentation or a temporary scar at the site of the sting
    \item Persisting itch for a few days
    \item A generalised allergic reaction, which is rare with bluebottles
    \item Delayed onset of more serious symptoms if the sting was actually from a different jellyfish
\end{itemize}

\section*{Prognosis and Outcomes}
Almost all bluebottle stings settle completely within a few hours to a few days, with pain easing shortly after hot water treatment. Serious or life-threatening reactions to bluebottle stings are extremely rare, in contrast to some other Australian jellyfish. The affected skin usually returns to normal within a week. If symptoms are severe, persist, or worsen, medical review is advisable.

\section*{Prevention and Self-Care}
\begin{itemize}
    \item Swim at patrolled beaches and follow lifeguard advice about jellyfish
    \item Avoid the water when bluebottles have been washed ashore or sighted in the surf
    \item Do not touch bluebottle jellyfish washed up on the beach -- their tentacles can still sting
    \item Wear protective clothing, such as a wetsuit or rash vest, when swimming in areas where jellyfish are common
    \item Carry a first aid kit with information on jellyfish stings when at remote beaches
    \item Learn the local jellyfish stinger first aid guidelines before heading to the water
\end{itemize}

\section*{Sources and Further Reading}
The following sources provide reliable, up-to-date information on bluebottle stings and marine first aid:

\begin{itemize}
    \item NHS. \textit{NHS guidance on jellyfish stings.} https://www.nhs.uk/conditions/
    \item Mayo Clinic. \textit{Information on jellyfish stings and first aid.} https://www.mayoclinic.org/
    \item MSD Manual. \textit{Overview of jellyfish stings and marine envenomation.} https://www.msdmanuals.com/
    \item Australian Red Cross. \textit{First aid guidance for marine stings.} https://www.redcross.org.au/
    \item MedlinePlus. \textit{Jellyfish stings patient information.} https://medlineplus.gov/
    \item World Health Organization. \textit{Injury prevention and first aid resources.} https://www.who.int/
\end{itemize}
